\chapter{Perspectives on the Modern Avatar of Caste}

\section{Backwardness and the Measurement Problem}

The central idea discussed (from a newspaper article the students were asked to read sociologically): \textbf{anybody can claim backwardness --- but how do you measure backwardness?} Communities now claim reservation by claiming backwardness.

\begin{KeyConcept}
\begin{itemize}
\item \textbf{SC} communities claim backwardness on the basis of \textbf{long social discrimination} with historical roots.
\item \textbf{ST} communities claim backwardness on the basis of \textbf{geographical isolation} (they were not mainstreamed; reservation was an \textbf{entry into the Western world / a gateway to meritocracy}, which was geographically isolated from the tribals).
\item \textbf{OBC} communities are '\textbf{other} backward' --- other than SC/ST --- and their backwardness is \textbf{social and educational}.
\item How do you \textbf{measure} social and educational backwardness? The \textbf{Mandal Commission} did this with its indicators/criteria --- but OBC status is \textbf{not} based on one uniform all-India scale: \textbf{every state has its own criteria for backwardness} (which makes OBC reservation the most contentious).
\item \textbf{EWS} (economically weaker sections) is a category, not a caste --- but it is \textbf{not open to those who already have reservation} (SC/ST/OBC are subtracted), so EWS means the economically weaker among the \textbf{upper castes/generals}.
\end{itemize}
\end{KeyConcept}

\begin{KeyConcept}
There are many communities that were once backward and have become \textbf{forward} (through reservation); and many that were once forward and have become \textbf{backward} --- backward relative to their own earlier position, or relative to those against whom they were once forward.
\end{KeyConcept}

\section{The Three Dimensions of Reservation Demands}

\begin{KeyConcept}
\begin{itemize}
\item \textbf{Dimension 1 --- reservation demand is made by dominant castes}, but the dominant caste is \textbf{not a homogeneous category}: the Yadavs are not dominant all over India (or even all over UP), and not all Yadavs are dominant. So who demands reservation? The \textbf{'backwards who are from the forwards'} --- the \textbf{backward forwards}.
\item \textbf{Dimension 2 --- the presence of 'backward forwards'} (backwards among the forwards) reflects \textbf{heterogeneity and division within the backwards}.
\item \textbf{Dimension 3 --- when the backward forwards demand reservation, it consolidates caste identities and weakens social justice}: they have class consciousness \emph{and} caste consciousness; they fight for economic upliftment but within the same caste --- so the class struggle carries a flavour of caste.
\item This yields a likely exam question: \textbf{'Reservation demands consolidate caste identity and weaken social justice. Comment.'}
\end{itemize}
\end{KeyConcept}

\begin{KeyConcept}
The backward forwards feel \textbf{relative deprivation} against the forward forwards, and want to \textbf{emulate} them --- earlier the lower castes emulated the upper castes; now there is \textbf{emulation within the caste} (the backward forwards emulating the forward forwards), which is a striving for \textbf{social mobility}.
\end{KeyConcept}

\section{Classization of Caste (D. L. Sheth)}

\begin{KeyConcept}
\begin{itemize}
\item The term \textbf{'classization of caste'} was given by the Indian sociologist \textbf{D. L. Sheth}.
\item It is a \textbf{two-fold process}:
\item 1. The \textbf{release of individual members of all castes from the religiously sanctioned status hierarchy} (the status hierarchy is no longer the determining factor).
\item 2. Their \textbf{interests and identities attach to organisations and categories} relevant to the \textbf{urban industrial system and modern politics} --- party membership is now more important than caste membership (as Beteille also said).
\item In short: \textbf{a process by which caste loses its traditional functions and acquires a new economic interest and political identity --- structural differentiation.}
\end{itemize}
\end{KeyConcept}

\begin{KeyConcept}
\begin{itemize}
\item Beteille's related point: \textbf{before modernization, caste was class; after modernization, there are classes within the caste} --- one caste can include a landlord, an officer and a cultivator (Brahmins, non-Brahmins and Adi Dravidas are all open to everything).
\item Caste is \textbf{evolving} --- like a snake shedding its old skin --- and this is the \textbf{new avatar of caste}.
\end{itemize}
\end{KeyConcept}

\section{AJGAR: The Dominant Castes}

\begin{KeyConcept}
\begin{itemize}
\item \textbf{M. N. Srinivas called the dominant caste an 'ajgar' (python)} --- the acronym \textbf{AJGAR = Ahir, Jat, Gujjar, Rajput} (the dominant castes; Ahir = Yadav).
\item These are the castes engaged in real-life reservation politics --- the 'backward forwards' demanding reservation.
\end{itemize}
\end{KeyConcept}

\begin{QuickRevision}
\begin{itemize}
\item Measurement of backwardness: SC (social discrimination), ST (geographical isolation; reservation = gateway to meritocracy), OBC (social + educational), EWS (economically weaker, but upper castes only).
\item Mandal Commission set OBC criteria, but every state has its own criteria of backwardness.
\item Three dimensions: (1) reservation demand by dominant castes (but dominant caste $\neq$ homogeneous); (2) backward forwards $\rightarrow$ heterogeneity within backwards; (3) reservation consolidates caste identity, weakens social justice.
\item Backward forwards feel relative deprivation, emulate forward forwards (emulation within caste = social mobility).
\item Classization of caste (D. L. Sheth): release from status hierarchy + attachment to modern organisations/politics; caste loses traditional functions, gains economic/political identity (structural differentiation).
\item Beteille: before modernization caste = class; after, classes within caste.
\item AJGAR (Srinivas's 'ajgar'/python) = Ahir, Jat, Gujjar, Rajput (dominant castes).
\end{itemize}
\end{QuickRevision}

